
\subsubsection{3.1 Problem Formulation}

Let M be a pre-trained decoder-only transformer with L layers. Standard LoRA fine-tuning trains low-rank adapter matrices {A_l, B_l} for each targeted layer l, with weight update ΔW_l = B_l A_l (rank r ≪ d). During training, the attention at each layer has access to the full KV cache: all T key-value pairs {(k_i, v_i)}_{i=1}^{T} participate in attention score computation.

At inference under KV cache budget ratio ρ ∈ (0, 1], only ⌊ρT⌋ KV pairs are retained. Under H2O eviction, the retained set consists of a fixed number of initial attention-sink tokens (n_sink) and the top-K tokens by cumulative attention score. Adapter parameters {A_l, B_l} trained with full attention are evaluated on a qualitatively different attention distribution.

\textbf{Goal:} Train {A_l, B_l} such that their representations are calibrated to the evicted-cache attention distribution encountered at inference under H2O eviction at budget ratio ρ.

\subsubsection{3.2 Eviction-Aware LoRA Training}

The approach applies H2O eviction masks during the LoRA forward pass at training time.

\#\#\#\# 3.2.1 H2O Mask Computation

For each training step, given input sequence of length T:

\begin{enumerate}
\item \textbf{Attention sink selection:} Retain the first n_sink tokens (default: 4) unconditionally.
\item \textbf{Heavy-hitter selection:} Compute cumulative attention scores s_i = Σ_{j≥i} α_{ji} where α_{ji} is the attention weight from position j to position i. Select the top-K positions by s_i, where K = ⌊ρT⌋ − n_sink.
\item \textbf{Mask construction:} M_i = 1 if position i is selected (sink or heavy hitter), 0 otherwise.
\end{enumerate}

The masked attention computation at layer l is:

    Attention(Q, K, V) = softmax((QK^T ⊙ M̃) / √d_k) V

where M̃_{ij} = M_j (column-wise mask; unselected positions set to −∞ before softmax).

\#\#\#\# 3.2.2 Hook-Based Mask Injection

The eviction mask is injected via a PyTorch forward hook on the attention module. This approach requires no architectural modifications and is compatible with HuggingFace \texttt{transformers} and \texttt{peft}. The hook intercepts the attention computation after Q, K, V projections have been computed but before the attention score matrix is finalized.

The hook is registered only during training forward passes and removed during evaluation and inference. This prevents double-application of eviction at inference and ensures the adapter is evaluated on the same mechanism it was trained to expect.

\textbf{Algorithm 1: Eviction-Aware LoRA Training}
```
Input:  Pre-trained model M, LoRA config (rank r, alpha α),
        training data D, KV budget ratio ρ, n_sink
Output: Adapter parameters {A_l, B_l}

\begin{enumerate}
\item Initialize LoRA adapters: B_l = 0, A_l ~ N(0, σ²) for all l
\item Register H2O forward hook on all attention layers:
\end{enumerate}
   hook(Q, K, V):
     if training:
       M = compute_h2o_mask(K, ρ, n_sink)
       return masked_attention(Q, K, V, M)
     else:
       return standard_attention(Q, K, V)
\begin{enumerate}
\item For each training batch (x, y) ∈ D:
\end{enumerate}
   a. Forward pass with hook active → eviction-masked attention
   b. Compute loss L(M(x), y) on LoRA parameters
   c. Backward pass → gradients reach {A_l, B_l} through eviction-masked activations
   d. Update {A_l, B_l} via AdamW
\begin{enumerate}
\item Deregister hook
\item Return {A_l, B_l}
\end{enumerate}
```

\#\#\#\# 3.2.3 Gradient Signal Analysis

Under eviction-aware training, ∂L/∂A_l flows through activations computed with masked attention: positions outside M contribute zero activation, and surviving positions carry the entire gradient signal. This differs from standard LoRA training, where ∂L/∂A_l flows through full-attention activations.

\subsubsection{3.3 Hyperparameters and Configuration}

\begin{table}[htbp]
\centering
\begin{tabular}{llll}
\toprule
Hyperparameter & Proxy Value & Target Value & Notes \\
\midrule
LoRA rank (r) & 8 & 16 & Rank 8 used in H-E1 smoke test \\
LoRA alpha (α) & 32 & 32 & Confirmed by experiment_results.json \\
LoRA dropout & 0.05 & 0.05 &  \\
KV budget ratio (ρ) & 0.50 & 0.50 & Primary comparison \\
Attention sinks (n_sink) & 4 & 4 & Standard H2O configuration \\
Training steps & 30 & ~1 epoch & Proxy: gradient-signal analysis only \\
Training data & LongAlpaca-12k (200 samples) & LongAlpaca-12k (full) &  \\
Optimizer & AdamW & AdamW & HuggingFace PEFT default \\
Learning rate & 2e-4 & 2e-4 &  \\
attn_implementation & eager & eager & Required for H2O+SDPA compatibility \\
\bottomrule
\end{tabular}
\end{table}

Note: The h-e1 validation report appendix recommends alpha=16 for dependent hypotheses (a forward-looking recommendation). The actual H-E1 training run used alpha=32, as confirmed by \texttt{experiment_results.json}.

\subsubsection{3.4 Relationship to Token-Scarcity Regularization}

Eviction-aware training can be interpreted as token-scarcity regularization — a training-time constraint that forces the adapter to develop representations under token-position absence. Key differences from dropout: (1) granularity — token-position eviction removes entire KV entries; dropout removes individual activations within each position; (2) determinism — H2O masks are policy-driven and input-dependent; dropout masks are stochastic and input-independent; (3) target — eviction-aware training addresses a specific training-inference distribution mismatch; dropout targets generalization.

\subsubsection{3.5 Infrastructure Components}

\textbf{H2O Training Wrapper:} Hook-based mask injection, validated for GPT-2 and architecturally validated for LLaMA-2-7B/Mistral-7B-v0.1.

\textbf{BudgetSweepEvaluator:} Evaluates adapters on LongBench (21 tasks) at ρ ∈ {0.25, 0.50, 0.75} in a single pass. Validated end-to-end at proxy scope (sequential GPT-2 adapter: 63 evaluations completed).

\textbf{SpearmanAnalyzer:} Computes Spearman ρ between KV budget ratio and per-category accuracy gap.

\textbf{AttentionEntropyExtractor:} Per-layer attention entropy and heavy-hitter concentration for paired statistical analysis.

